% Options for packages loaded elsewhere
% Options for packages loaded elsewhere
\PassOptionsToPackage{unicode}{hyperref}
\PassOptionsToPackage{hyphens}{url}
\PassOptionsToPackage{dvipsnames,svgnames,x11names}{xcolor}
%
\documentclass[
  letterpaper,
  DIV=11,
  numbers=noendperiod]{scrartcl}
\usepackage{xcolor}
\usepackage{amsmath,amssymb}
\setcounter{secnumdepth}{-\maxdimen} % remove section numbering
\usepackage{iftex}
\ifPDFTeX
  \usepackage[T1]{fontenc}
  \usepackage[utf8]{inputenc}
  \usepackage{textcomp} % provide euro and other symbols
\else % if luatex or xetex
  \usepackage{unicode-math} % this also loads fontspec
  \defaultfontfeatures{Scale=MatchLowercase}
  \defaultfontfeatures[\rmfamily]{Ligatures=TeX,Scale=1}
\fi
\usepackage{lmodern}
\ifPDFTeX\else
  % xetex/luatex font selection
\fi
% Use upquote if available, for straight quotes in verbatim environments
\IfFileExists{upquote.sty}{\usepackage{upquote}}{}
\IfFileExists{microtype.sty}{% use microtype if available
  \usepackage[]{microtype}
  \UseMicrotypeSet[protrusion]{basicmath} % disable protrusion for tt fonts
}{}
\makeatletter
\@ifundefined{KOMAClassName}{% if non-KOMA class
  \IfFileExists{parskip.sty}{%
    \usepackage{parskip}
  }{% else
    \setlength{\parindent}{0pt}
    \setlength{\parskip}{6pt plus 2pt minus 1pt}}
}{% if KOMA class
  \KOMAoptions{parskip=half}}
\makeatother
% Make \paragraph and \subparagraph free-standing
\makeatletter
\ifx\paragraph\undefined\else
  \let\oldparagraph\paragraph
  \renewcommand{\paragraph}{
    \@ifstar
      \xxxParagraphStar
      \xxxParagraphNoStar
  }
  \newcommand{\xxxParagraphStar}[1]{\oldparagraph*{#1}\mbox{}}
  \newcommand{\xxxParagraphNoStar}[1]{\oldparagraph{#1}\mbox{}}
\fi
\ifx\subparagraph\undefined\else
  \let\oldsubparagraph\subparagraph
  \renewcommand{\subparagraph}{
    \@ifstar
      \xxxSubParagraphStar
      \xxxSubParagraphNoStar
  }
  \newcommand{\xxxSubParagraphStar}[1]{\oldsubparagraph*{#1}\mbox{}}
  \newcommand{\xxxSubParagraphNoStar}[1]{\oldsubparagraph{#1}\mbox{}}
\fi
\makeatother


\usepackage{longtable,booktabs,array}
\usepackage{calc} % for calculating minipage widths
% Correct order of tables after \paragraph or \subparagraph
\usepackage{etoolbox}
\makeatletter
\patchcmd\longtable{\par}{\if@noskipsec\mbox{}\fi\par}{}{}
\makeatother
% Allow footnotes in longtable head/foot
\IfFileExists{footnotehyper.sty}{\usepackage{footnotehyper}}{\usepackage{footnote}}
\makesavenoteenv{longtable}
\usepackage{graphicx}
\makeatletter
\newsavebox\pandoc@box
\newcommand*\pandocbounded[1]{% scales image to fit in text height/width
  \sbox\pandoc@box{#1}%
  \Gscale@div\@tempa{\textheight}{\dimexpr\ht\pandoc@box+\dp\pandoc@box\relax}%
  \Gscale@div\@tempb{\linewidth}{\wd\pandoc@box}%
  \ifdim\@tempb\p@<\@tempa\p@\let\@tempa\@tempb\fi% select the smaller of both
  \ifdim\@tempa\p@<\p@\scalebox{\@tempa}{\usebox\pandoc@box}%
  \else\usebox{\pandoc@box}%
  \fi%
}
% Set default figure placement to htbp
\def\fps@figure{htbp}
\makeatother


% definitions for citeproc citations
\NewDocumentCommand\citeproctext{}{}
\NewDocumentCommand\citeproc{mm}{%
  \begingroup\def\citeproctext{#2}\cite{#1}\endgroup}
\makeatletter
 % allow citations to break across lines
 \let\@cite@ofmt\@firstofone
 % avoid brackets around text for \cite:
 \def\@biblabel#1{}
 \def\@cite#1#2{{#1\if@tempswa , #2\fi}}
\makeatother
\newlength{\cslhangindent}
\setlength{\cslhangindent}{1.5em}
\newlength{\csllabelwidth}
\setlength{\csllabelwidth}{3em}
\newenvironment{CSLReferences}[2] % #1 hanging-indent, #2 entry-spacing
 {\begin{list}{}{%
  \setlength{\itemindent}{0pt}
  \setlength{\leftmargin}{0pt}
  \setlength{\parsep}{0pt}
  % turn on hanging indent if param 1 is 1
  \ifodd #1
   \setlength{\leftmargin}{\cslhangindent}
   \setlength{\itemindent}{-1\cslhangindent}
  \fi
  % set entry spacing
  \setlength{\itemsep}{#2\baselineskip}}}
 {\end{list}}
\usepackage{calc}
\newcommand{\CSLBlock}[1]{\hfill\break\parbox[t]{\linewidth}{\strut\ignorespaces#1\strut}}
\newcommand{\CSLLeftMargin}[1]{\parbox[t]{\csllabelwidth}{\strut#1\strut}}
\newcommand{\CSLRightInline}[1]{\parbox[t]{\linewidth - \csllabelwidth}{\strut#1\strut}}
\newcommand{\CSLIndent}[1]{\hspace{\cslhangindent}#1}



\setlength{\emergencystretch}{3em} % prevent overfull lines

\providecommand{\tightlist}{%
  \setlength{\itemsep}{0pt}\setlength{\parskip}{0pt}}



 


\usepackage{hyperref}
\KOMAoption{captions}{tableheading}
\makeatletter
\@ifpackageloaded{tcolorbox}{}{\usepackage[skins,breakable]{tcolorbox}}
\@ifpackageloaded{fontawesome5}{}{\usepackage{fontawesome5}}
\definecolor{quarto-callout-color}{HTML}{909090}
\definecolor{quarto-callout-note-color}{HTML}{0758E5}
\definecolor{quarto-callout-important-color}{HTML}{CC1914}
\definecolor{quarto-callout-warning-color}{HTML}{EB9113}
\definecolor{quarto-callout-tip-color}{HTML}{00A047}
\definecolor{quarto-callout-caution-color}{HTML}{FC5300}
\definecolor{quarto-callout-color-frame}{HTML}{acacac}
\definecolor{quarto-callout-note-color-frame}{HTML}{4582ec}
\definecolor{quarto-callout-important-color-frame}{HTML}{d9534f}
\definecolor{quarto-callout-warning-color-frame}{HTML}{f0ad4e}
\definecolor{quarto-callout-tip-color-frame}{HTML}{02b875}
\definecolor{quarto-callout-caution-color-frame}{HTML}{fd7e14}
\makeatother
\makeatletter
\@ifpackageloaded{caption}{}{\usepackage{caption}}
\AtBeginDocument{%
\ifdefined\contentsname
  \renewcommand*\contentsname{Table of contents}
\else
  \newcommand\contentsname{Table of contents}
\fi
\ifdefined\listfigurename
  \renewcommand*\listfigurename{List of Figures}
\else
  \newcommand\listfigurename{List of Figures}
\fi
\ifdefined\listtablename
  \renewcommand*\listtablename{List of Tables}
\else
  \newcommand\listtablename{List of Tables}
\fi
\ifdefined\figurename
  \renewcommand*\figurename{Figure}
\else
  \newcommand\figurename{Figure}
\fi
\ifdefined\tablename
  \renewcommand*\tablename{Table}
\else
  \newcommand\tablename{Table}
\fi
}
\@ifpackageloaded{float}{}{\usepackage{float}}
\floatstyle{ruled}
\@ifundefined{c@chapter}{\newfloat{codelisting}{h}{lop}}{\newfloat{codelisting}{h}{lop}[chapter]}
\floatname{codelisting}{Listing}
\newcommand*\listoflistings{\listof{codelisting}{List of Listings}}
\makeatother
\makeatletter
\makeatother
\makeatletter
\@ifpackageloaded{caption}{}{\usepackage{caption}}
\@ifpackageloaded{subcaption}{}{\usepackage{subcaption}}
\makeatother
\usepackage{bookmark}
\IfFileExists{xurl.sty}{\usepackage{xurl}}{} % add URL line breaks if available
\urlstyle{same}
\hypersetup{
  pdftitle={Interactive R tutorials for learning and assessments},
  pdfauthor={Léo Belzile; David A. Campbell; Kyran Cupido; Nishan Mudalige; Aurélien Nicosia; Emily Somerset; Haochen Song; David Riegert; Russell Steele; Irene Vrbik; Karen Huynh Wong},
  pdfkeywords={R programming language; automated feedback;},
  colorlinks=true,
  linkcolor={blue},
  filecolor={Maroon},
  citecolor={Blue},
  urlcolor={Blue},
  pdfcreator={LaTeX via pandoc}}


\title{Interactive R tutorials for learning and assessments}
\author{Léo
Belzile\thanks{HEC Montréal. Email: \href{mailto:leo.belzile@hec.ca}{leo.belzile@hec.ca}. ORCID: \href{https://orcid.org/0000-0002-9135-014X}{0000-0002-9135-014X}.} \and David
A.
Campbell\thanks{Carleton University. Email: \href{mailto:davecampbell@math.carleton.ca}{davecampbell@math.carleton.ca}. ORCID: \href{https://orcid.org/0000-0001-8722-4450}{0000-0001-8722-4450}.} \and Kyran
Cupido\thanks{St. Francis Xavier University. Email: \href{mailto:kcupido@stfx.ca}{kcupido@stfx.ca}. ORCID: \href{https://orcid.org/0000-0001-6984-4082}{0000-0001-6984-4082}.} \and Nishan
Mudalige\thanks{University of Toronto. Email: \href{mailto:nishan.mudalige@utoronto.ca}{nishan.mudalige@utoronto.ca}. ORCID: \href{https://orcid.org/0000-0003-3981-9868}{0000-0003-3981-9868}.} \and Aurélien
Nicosia\thanks{Université Laval. Email: \href{mailto:aurelien.nicosia@mat.ulaval.ca}{aurelien.nicosia@mat.ulaval.ca}.} \and Emily
Somerset\thanks{University of Toronto. Email: \href{mailto:emily.somerset@mail.utoronto.ca}{emily.somerset@mail.utoronto.ca}.} \and Haochen
Song\thanks{University of Toronto. Email: \href{mailto:fred.song@mail.utoronto.ca}{fred.song@mail.utoronto.ca}. ORCID: \href{https://orcid.org/0009-0006-7825-4993}{0009-0006-7825-4993}.} \and David
Riegert\thanks{Trent University. Email: \href{mailto:davidriegert@trentu.ca}{davidriegert@trentu.ca}.} \and Russell
Steele\thanks{McGill University. Email: \href{mailto:russell.steele@mcgill.ca}{russell.steele@mcgill.ca}. ORCID: \href{https://orcid.org/0000-0003-1071-0914}{0000-0003-1071-0914}.} \and Irene
Vrbik\thanks{University of British Columbia. Email: \href{mailto:irene.vrbik@ubc.ca}{irene.vrbik@ubc.ca}.} \and Karen
Huynh
Wong\thanks{University of Toronto. Email: \href{mailto:karen.huynhwong@utoronto.ca}{karen.huynhwong@utoronto.ca}.}}
\date{}
\begin{document}
\maketitle
\begin{abstract}
We review software tools for teaching, assessing and grading data
analysis skills in the classroom using the \textbf{R} programming
language.
\end{abstract}


\subsection{Introduction}\label{introduction}

Data science majors popularity is at an all time high (Pierson 2023),
leading to strain on teaching resources. The increased enrollment, even
in upper level undergraduate courses, increases the time spent on class
management by instructors and limits their ability and that of teaching
assistants to provide timely feedback on evaluations. The advent of
generative artificial intelligence (genAI) also raises concerns about
plagiarism. Due to these, instructors more and more default back to
invigilated multiple choice exams. This is unfortunate because the
questions that can be asked in such evaluations are not necessarily
authentic.

Data science is interdisciplinary by nature: standard tasks performed by
analysts include exploratory data analysis and data wrangling,
visualisation, model specification, estimation and validation, along
with reporting. Designing assignments should therefore try to emulate
those tasks while fostering statistical thinking (Woodard, Lee, and
Woodard 2020). For formative and summative assessment, it is often
beneficial to have clear and consistent grading rubrics, which
facilitate learning experience (Poulos and Mahony 2008). Automating
feedback by parsing code, and using natural language processing (NLP)
for comments can lead to efficiency gains (Galassi and Vittorini 2021),
while software solutions such as, e.g., Otter-Grader (Pyles and UC
Berkeley Data Science Education Program 2025) facilitate sharing of
grading rubrics, ensuring some consistency across instructors. Version
control is now taught early on in introductory data science courses
(e.g., Timbers, Campbell, and Lee 2022), and it is also becoming more
frequent for class management (Rundel and Cetinkaya-Rundel 2025; Ricci,
Medina, and Dogucu 2024). Solutions such as WebWork (Jahodova Berkova
and Nemec 2020; Zimmerman 2020) or R-exams (Zeileis, Umlauf, and Leisch
2014) gives instructors tools to effectively manage large-scale
assessments using randomized questions databases. Automation of many of
these grading and feedback tasks remains difficult because, while there
are wrong answers, there is no unique correct one. Thus, setting up
rubrics requires much upfront work to determine most likely errors.

Our purpose is to review software tools for teaching, assessing and
grading data analysis skills in the classroom using the \textbf{R}
programming language (R Core Team 2025), a widely-used interpreted
language which is popular among statisticians. There exists a wealth of
contributed packages and software bundles to encourage semi-automation
of tasks, generation of randomized evaluations, and tutorials. One key
point is keeping the technical barriers low also for instructors, who
can easily adapt their existing examples to a more interactive format.
We consider several examples, notably:

\begin{itemize}
\tightlist
\item
  in-class activities with limited entry-level barriers to programming
  via \href{https://docs.r-wasm.org/webr/latest/}{\texttt{webR}} coupled
  with \href{https://r-wasm.github.io/quarto-live/}{Quarto Live};
\item
  interactive tutorials via the \texttt{learnR} (Aden-Buie et al. 2025;
  Stoudt, Scotina, and Luebke 2022a) package, including description of
  randomization within, with gradated hints and support for multiple
  solutions
  \href{https://pkgs.rstudio.com/gradethis/}{\texttt{gradethis}} and
  recovery of student responses using
  \href{https://dtkaplan.github.io/submitr/articles/using.html}{\texttt{submitr}}
  or \href{https://github.com/rundel/learnrhash}{\texttt{learnhash}}.
\item
  semiautomatic tools for extracting student solutions for coding tasks
  and autograding, focusing on invariance of outputs.
\end{itemize}

The purpose of this paper is to review these three aspects critically
with different points in mind, and to

\begin{itemize}
\tightlist
\item
  streamline use of these tools, by providing templates for users,
\item
  review and highlight functionalities (e.g., limitations in
  \texttt{learnr} for randomization due to pre-rendering),
\item
  suggest strategies for building unit tests and enabling informative
  feedback,
\item
  report on instructors experience implementing these tools in their
  classroom during the 2025--2026 academic year.
\end{itemize}

\subsection{WebR interface and interactive online
tutorials}\label{webr-interface-and-interactive-online-tutorials}

\subsection{Description of tools}\label{description-of-tools}

\begin{longtable}[]{@{}
  >{\raggedright\arraybackslash}p{(\linewidth - 6\tabcolsep) * \real{0.1647}}
  >{\raggedright\arraybackslash}p{(\linewidth - 6\tabcolsep) * \real{0.2706}}
  >{\raggedright\arraybackslash}p{(\linewidth - 6\tabcolsep) * \real{0.2706}}
  >{\raggedright\arraybackslash}p{(\linewidth - 6\tabcolsep) * \real{0.2706}}@{}}
\toprule\noalign{}
\begin{minipage}[b]{\linewidth}\raggedright
Tool
\end{minipage} & \begin{minipage}[b]{\linewidth}\raggedright
Description
\end{minipage} & \begin{minipage}[b]{\linewidth}\raggedright
Hosting
\end{minipage} & \begin{minipage}[b]{\linewidth}\raggedright
Related packages
\end{minipage} \\
\midrule\noalign{}
\endhead
\bottomrule\noalign{}
\endlastfoot
\texttt{webR} & A version of R compiled for the browser and Node.js
using WebAssembly. It allows R code to run in the user's browser without
an R server (r-wasm contributors 2026b). & Can work with static web
hosting because computation happens client-side, although package
support depends on WebAssembly- compatible R packages (r-wasm
contributors 2026a, 2026b). & Useful for lightweight, browser-based R
interactivity. It is conceptually different from standard
\texttt{learnr}, which usually depends on Shiny. \\
\texttt{learnr} & An R package for turning R Markdown documents into
interactive tutorials with narrative, code exercises, quizzes, videos,
and Shiny components (Aden-Buie et al. 2025). & Standard \texttt{learnr}
tutorials use \texttt{runtime:\ shiny\_prerendered}, so they usually
need to run locally or on a Shiny-capable hosting platform rather than
as ordinary static pages (Aden-Buie et al. 2025). & Provides the
tutorial framework that can incorporate \texttt{gradethis} feedback and,
in some workflows, \texttt{learnrhash} tracking. \\
\texttt{gradethis} & An R package that provides automated feedback for
student exercises in \texttt{learnr} tutorials. It can compare student
code to a model solution or use custom grading logic (Aden-Buie et al.
2023). & Its deployment constraints generally follow \texttt{learnr},
since it is designed to run inside \texttt{learnr} tutorials (Aden-Buie
et al. 2023). & Extends \texttt{learnr} by adding automated feedback to
code exercises. \\
\texttt{Quarto} & An open-source scientific and technical publishing
system for documents, websites, blogs, books, presentations, dashboards,
and other outputs (Posit Software, PBC 2026). & Can produce static HTML
outputs that are suitable for static hosting; dynamic or live
interactive content may require an appropriate runtime or hosting
environment (Posit Software, PBC 2026). & Useful for publishing course
materials, while \texttt{learnr} is more specialized for interactive
tutorials. \\
\texttt{Other\ tools?} & & & \\
\end{longtable}

\subsection{User case}\label{user-case}

\begin{longtable}[]{@{}
  >{\raggedright\arraybackslash}p{(\linewidth - 6\tabcolsep) * \real{0.1647}}
  >{\raggedright\arraybackslash}p{(\linewidth - 6\tabcolsep) * \real{0.2706}}
  >{\raggedright\arraybackslash}p{(\linewidth - 6\tabcolsep) * \real{0.2706}}
  >{\raggedright\arraybackslash}p{(\linewidth - 6\tabcolsep) * \real{0.2706}}@{}}
\toprule\noalign{}
\begin{minipage}[b]{\linewidth}\raggedright
Tool
\end{minipage} & \begin{minipage}[b]{\linewidth}\raggedright
Use case summary
\end{minipage} & \begin{minipage}[b]{\linewidth}\raggedright
Strengths
\end{minipage} & \begin{minipage}[b]{\linewidth}\raggedright
Limitations / comparison
\end{minipage} \\
\midrule\noalign{}
\endhead
\bottomrule\noalign{}
\endlastfoot
\texttt{webR} & Use when students need to run R code directly in the
browser without installing R or connecting to an R server (r-wasm
contributors 2026b). & Strong for static or mostly-static websites
because computation happens client-side. Useful for lightweight
browser-based coding activities, demos, and interactive examples. & Less
mature than standard server-side R workflows. The API is under active
development, some packages may not work, and mobile browsers may have
memory limits (r-wasm contributors 2026b). Compared with
\texttt{learnr}, it is better for static hosting but less specialized
for structured tutorials and feedback. \\
\texttt{learnr} & Use when you want a guided interactive tutorial with
narrative, editable code exercises, quizzes, videos, and Shiny
components (Aden-Buie et al. 2025). & Strong for self-paced active
learning. Gives students a structured environment where they read, run
code, answer questions, and receive feedback. Works well with
\texttt{gradethis}. & Standard \texttt{learnr} tutorials use
\texttt{runtime:\ shiny\_prerendered}, so they usually need a
Shiny-capable environment rather than ordinary static hosting (Aden-Buie
et al. 2025). Compared with Quarto, it is more specialized for tutorials
but less general for course websites, slides, books, or reports. \\
\texttt{gradethis} & Use when you want automated feedback on student
code inside a \texttt{learnr} tutorial (Aden-Buie et al. 2023). & Strong
for formative assessment. Can compare student code to a model solution
or use custom grading logic, which helps give targeted feedback for
common mistakes (Aden-Buie et al. 2023). & It is not a general
publishing tool or a standalone assessment system. Its use is closely
tied to \texttt{learnr}, so its hosting constraints usually follow
\texttt{learnr}. Compared with \texttt{learnr}, it is not the tutorial
framework itself; it is the feedback layer. \\
\texttt{Quarto} & Use when you want to create course notes, reports,
slides, websites, blogs, books, dashboards, or other reproducible
teaching materials (Posit Software, PBC 2026). & Strong for publishing.
Supports multiple languages, including R, Python, Julia, and Observable,
and can produce many output formats such as HTML, PDF, Word, slides,
websites, and books (Posit Software, PBC 2026). & It is not primarily an
interactive tutorial or grading system. Static outputs are easy to host,
but live computation or rich interactivity may require additional tools
or hosting. Compared with \texttt{learnr}, Quarto is broader and better
for publishing; \texttt{learnr} is better for guided coding
tutorials. \\
\end{longtable}

\subsection{\texorpdfstring{Interactive tutorials for learning and
assessment using
\texttt{learnR}}{Interactive tutorials for learning and assessment using learnR}}\label{interactive-tutorials-for-learning-and-assessment-using-learnr}

\subsubsection{Main initial takeaways}\label{main-initial-takeaways}

\begin{itemize}
\tightlist
\item
  Best use is for interactive, low stakes tutorials.
\item
  Supports active learning since students learn by running code.
\item
  Good for formative assessment: useful for practice, readiness checks,
  lab prep, and review rather than formal high-stakes grading.
\item
  Works best for smaller code tasks to simplify feedback through
  \texttt{gradethis}.
\item
  Immediate feedback is available with \texttt{gradethis}.
\item
  Student work can be tracked easily through \texttt{learnrhash}, as
  long as you need an existence proof and not a formal evaluation.
\item
  The choice of hosting matters. Static tutorial \texttt{learnr}
  tutorials materials or pre-rendered R Markdown pages can be shared
  through GitHub Pages, but standard live interactive \texttt{learnr}
  tutorials require a Shiny-capable environment.
\end{itemize}

Some issues to think about:

\begin{itemize}
\tightlist
\item
  There are scoping issues which require careful planning when designing
  tutorials.
\item
  Security is potentially an issue with regards to what you can do,
  particularly if you are serving the learnr tutorial from a server.
\item
  It requires work to convert output from \texttt{learnr} hash into
  readable content for more formal grading.
\item
  For questions where the next part of a question depends on the
  previous part, you should give the student the answer.
\item
  \texttt{learnrhash} relies on Shiny server logic to collect the
  tutorial state and generate the copy-paste hash (\texttt{learnrhash}
  can not be implemented on GitHub pages as far as I am aware)
\end{itemize}

\subsection{Autograding of code chunks from Rmarkdown and Quarto
documents}\label{autograding-of-code-chunks-from-rmarkdown-and-quarto-documents}

\subsection{Designing tutorials and
assessment}\label{designing-tutorials-and-assessment}

In the following section, we focus on some elements that are
idiosyncratic to \textbf{R}, although much of the discussion has similar
implications for any other programming language.

\begin{itemize}
\tightlist
\item
  A simple way of getting around the cheating problem is through
  individualized data generation for which the procedure/script remains
  the same but answers differ. Of course, the drawback of this choice is
  that scripts could be shared between students. In data analysis
  problems, it is also possible to randomize the order of covariates,
  their name and their selection.
\item
  Asking the right questions: determining the variables and setup (e.g.,
  test statistic, one-sided versus two-sided test, confidence level,
  conclusion) that can be randomized.
\item
  Generating context using generative AI.
\item
  Validating the results (e.g., numerical tolerance to rounding,
  differences for p-values for one-sided versus two-sided).
\end{itemize}

Summative assessment of student answers

\begin{itemize}
\tightlist
\item
  Open-ended questions are best checked through multiple choice
  questions that force users to generate plots or look at summary
  statistics and model output.
\item
  Focus on invariants and account for students degrees of freedom: e.g.,
  in regression models, we can't check formulas because variable names
  may differ following data wrangling, the syntax and order of variables
  may be different (in \textbf{R} formula notation, \texttt{a*b} versus
  \texttt{b\ +\ a\ +\ a:b}). Likewise, contrast specification may
  change, as is the default engine class (e.g., \texttt{lm} vs
  \texttt{glm} with \texttt{family=gaussian}) for fitting these models.
  We can however compare sample sizes, span of projection (hat)
  matrices, linear contrasts, tests comparing nested models, residuals
  and predictions, etc.
\item
  Depending on how prescriptive instructors are (for example, with
  naming of variables during tidying setup), we can validate more
  output.
\end{itemize}

Students could simply save certain objects as lists and submit a .RData
or a CSV file with their short answer. Automated checks can verify class
or attributes, and check existence of the variable names; settings may
also be set so that the environment is saved at the end of each
question. This requires that the answers be numerical, or TRUE/FALSE. If
the inputs are randomized, the instructor must run every data version to
compare student answers with the solution key.

If students submit a Quarto file that generates an html/PDF report.
consider chopping a Quarto file with headers defining questions and
sub-questions to facilitate extraction of code or run a text parser to
extract comments or text.

\subsection{Evidence from the
classroom}\label{evidence-from-the-classroom}

\begin{tcolorbox}[enhanced jigsaw, bottomrule=.15mm, title=\textcolor{quarto-callout-caution-color}{\faFire}\hspace{0.5em}{Caution}, opacityback=0, left=2mm, colbacktitle=quarto-callout-caution-color!10!white, toprule=.15mm, breakable, titlerule=0mm, colback=white, coltitle=black, leftrule=.75mm, opacitybacktitle=0.6, colframe=quarto-callout-caution-color-frame, arc=.35mm, bottomtitle=1mm, rightrule=.15mm, toptitle=1mm]

The list below has not been thoroughly checked and is not exhaustive.

\end{tcolorbox}

\subsubsection{\texorpdfstring{\texttt{learnr}}{learnr}}\label{learnr}

\begin{longtable}[]{@{}
  >{\raggedright\arraybackslash}p{(\linewidth - 6\tabcolsep) * \real{0.2299}}
  >{\raggedright\arraybackslash}p{(\linewidth - 6\tabcolsep) * \real{0.1897}}
  >{\raggedright\arraybackslash}p{(\linewidth - 6\tabcolsep) * \real{0.2241}}
  >{\raggedright\arraybackslash}p{(\linewidth - 6\tabcolsep) * \real{0.3448}}@{}}
\toprule\noalign{}
\begin{minipage}[b]{\linewidth}\raggedright
Source
\end{minipage} & \begin{minipage}[b]{\linewidth}\raggedright
Evidence type
\end{minipage} & \begin{minipage}[b]{\linewidth}\raggedright
Classroom / teaching context
\end{minipage} & \begin{minipage}[b]{\linewidth}\raggedright
Summary
\end{minipage} \\
\midrule\noalign{}
\endhead
\bottomrule\noalign{}
\endlastfoot
Stoudt, Scotina, \& Luebke, ``Supporting Statistics and Data Science
Education with learnr'' (Stoudt, Scotina, and Luebke 2022b) &
Peer-reviewed education article & Statistics and data science education
& Discusses classroom uses of \texttt{learnr} tutorials for statistics
and data science instruction. \\
Engels, Grosjean, \& Artus, ``Teaching data science to students in
biology using R, RStudio and Learnr'' (Engels, Grosjean, and Artus 2023)
& Peer-reviewed classroom study & Biology students learning data science
& Analyzes three years of course data and describes \texttt{learnr}
tutorials as part of the course design. \\
Aberson, ``Building Interactive Tutorials for Teaching Psychological
Statistics Online with learnr'' (Aberson 2021) & Peer-reviewed teaching
article & Psychological statistics & Describes building online
interactive tutorials with videos, quizzes, and R analysis space using
\texttt{learnr}. \\
Stoudt, ``Using learnr Tutorials in an Intro Stats Class as Pre-Labs''
(Stoudt 2021) & Instructor classroom report / blog & Introductory
statistics pre-labs & Practical classroom account of using
\texttt{learnr} tutorials before lab sessions. Not peer-reviewed. \\
\end{longtable}

\subsubsection{\texorpdfstring{\texttt{webR}}{webR}}\label{webr}

\begin{longtable}[]{@{}
  >{\raggedright\arraybackslash}p{(\linewidth - 6\tabcolsep) * \real{0.2299}}
  >{\raggedright\arraybackslash}p{(\linewidth - 6\tabcolsep) * \real{0.1897}}
  >{\raggedright\arraybackslash}p{(\linewidth - 6\tabcolsep) * \real{0.2241}}
  >{\raggedright\arraybackslash}p{(\linewidth - 6\tabcolsep) * \real{0.3448}}@{}}
\toprule\noalign{}
\begin{minipage}[b]{\linewidth}\raggedright
Source
\end{minipage} & \begin{minipage}[b]{\linewidth}\raggedright
Evidence type
\end{minipage} & \begin{minipage}[b]{\linewidth}\raggedright
Classroom / teaching context
\end{minipage} & \begin{minipage}[b]{\linewidth}\raggedright
Summary
\end{minipage} \\
\midrule\noalign{}
\endhead
\bottomrule\noalign{}
\endlastfoot
Perkel, ``No installation required: how WebAssembly is changing
scientific computing'' (Perkel 2024) & Nature technology feature &
Education-related use case & Discusses WebAssembly scientific computing
and includes browser-based R as a way to reduce installation barriers.
Not a classroom-outcomes study. \\
Rennie, ``Teaching statistics interactively with webR'' (Rennie 2023) &
Conference talk / teaching demonstration & Statistics teaching &
Demonstrates how \texttt{webR} can be used to make teaching materials
more interactive with live coding examples. Not peer-reviewed. \\
White, ``A brief introduction to using WebR and Quarto for client-side
interactive lesson material'' (White 2023) & Teaching blog / technical
demonstration & Client-side interactive lesson materials & Shows how
\texttt{webR} and Quarto can be used to let students run R code in
lesson materials without installing R. Not peer-reviewed. \\
Ward, ``stan-playground: Run Stan models directly in your browser''
(Ward, Soules, and Magland 2026) & Peer-reviewed software paper &
Browser-based statistical modeling for instruction & Not a classroom
study of \texttt{webR}, but relevant to browser-based statistical
computing for instruction; discusses integration with \texttt{webR}. \\
\end{longtable}

\phantomsection\label{refs}
\begin{CSLReferences}{1}{0}
\bibitem[\citeproctext]{ref-aberson2021learnrpsych}
Aberson, Christopher L. 2021. {``Building Interactive Tutorials for
Teaching Psychological Statistics Online with Learnr.''}
\emph{Technology Innovations in Statistics Education} 13 (1).
\url{https://doi.org/10.5070/T513153822}.

\bibitem[\citeproctext]{ref-gradethis}
Aden-Buie, Garrick, Daniel Chen, Garrett Grolemund, Alexander Rossell
Hayes, and Barret Schloerke. 2023. \emph{{gradethis}: Automated Feedback
for Student Exercises in {learnr} Tutorials}. Posit Software, PBC.
\url{https://pkgs.rstudio.com/gradethis/}.

\bibitem[\citeproctext]{ref-learnr}
Aden-Buie, Garrick, Barret Schloerke, JJ Allaire, and Alexander Rossell
Hayes. 2025. \emph{{learnr}: Interactive Tutorials for {R}}. Posit
Software, PBC. \url{https://doi.org/10.32614/CRAN.package.learnr}.

\bibitem[\citeproctext]{ref-engels2023learnrbiology}
Engels, Guyliann, Philippe Grosjean, and Frédérique Artus. 2023.
{``Teaching Data Science to Students in Biology Using r, RStudio and
Learnr: Analysis of Three Years Data.''} \emph{Foundations of Data
Science} 5 (2): 266--85. \url{https://doi.org/10.3934/fods.2022022}.

\bibitem[\citeproctext]{ref-Galassi.Vittorini:2021}
Galassi, Alessandra, and Pierpaolo Vittorini. 2021. {``Automated
Feedback to Students in Data Science Assignments: Improved
Implementation and Results.''} In \emph{Proceedings of the 14th Biannual
Conference of the Italian SIGCHI Chapter}, 8. CHItaly '21 3. New York,
NY, USA: Association for Computing Machinery.
\url{https://doi.org/10.1145/3464385.3464387}.

\bibitem[\citeproctext]{ref-Berkova.Nemec:2020}
Jahodova Berkova, Andrea, and Radek Nemec. 2020. {``Teaching Theory of
Probability and Statistics During the {C}ovid-19 Emergency.''}
\emph{Symmetry} 12 (9): 1577. \url{https://doi.org/10.3390/sym12091577}.

\bibitem[\citeproctext]{ref-perkel2024webassembly}
Perkel, Jeffrey M. 2024. {``No Installation Required: How WebAssembly Is
Changing Scientific Computing.''} \emph{Nature}.
\url{https://doi.org/10.1038/d41586-024-00725-1}.

\bibitem[\citeproctext]{ref-Pierson:2023}
Pierson, Steve. 2023. {``Data Analytics, Data Science Degrees See Large
Increases in 2022.''} \emph{AMStat News}, no. 558: 16--21.
\url{https://magazine.amstat.org/blog/2023/12/01/degreesstats2022/}.

\bibitem[\citeproctext]{ref-quarto}
Posit Software, PBC. 2026. \emph{{Quarto}: An Open-Source Scientific and
Technical Publishing System}. Posit Software, PBC.
\url{https://quarto.org/}.

\bibitem[\citeproctext]{ref-Poulos.Mahony:2008}
Poulos, Ann, and Mary Jane Mahony. 2008. {``Effectiveness of Feedback:
The Students' Perspective.''} \emph{Assessment \& Evaluation in Higher
Education} 33 (2): 143--54.
\url{https://doi.org/10.1080/02602930601127869}.

\bibitem[\citeproctext]{ref-ottergrader:2025}
Pyles, Christopher, and UC Berkeley Data Science Education Program.
2025. {``{Otter-Grader}: A {P}ython and {R} Autograding Solution.''}
\url{https://doi.org/10.5281/zenodo.5259955}.

\bibitem[\citeproctext]{ref-Rlang}
R Core Team. 2025. \emph{R: A Language and Environment for Statistical
Computing}. Vienna, Austria: R Foundation for Statistical Computing.
\url{https://www.R-project.org/}.

\bibitem[\citeproctext]{ref-rennie2023webr}
Rennie, Nicola. 2023. {``Teaching Statistics Interactively with webR.''}
Talk at the RSS International Conference 2023.
\url{https://nrennie.rbind.io/talks/rss-conference-teaching-webr/}.

\bibitem[\citeproctext]{ref-Ricci.Medina.Dogucu:2024}
Ricci, Federica Zoe, Catalina Mari Medina, and Mine Dogucu. 2024.
{``Automated Grading Workflows for Providing Personalized Feedback to
Open-Ended Data Science Assignments.''} \emph{Technology Innovations in
Statistics Education} 15 (1). \url{https://doi.org/10.5070/t5.1886}.

\bibitem[\citeproctext]{ref-ghclass}
Rundel, Colin, and Mine Cetinkaya-Rundel. 2025. \emph{{ghclass}: Tools
for Managing Classes on {G}itHub}.
\url{https://github.com/rundel/ghclass}.

\bibitem[\citeproctext]{ref-webRpackages}
r-wasm contributors. 2026a. \emph{Installing {R} Packages for {webR}}.
{r-wasm}. \url{https://docs.r-wasm.org/webr/latest/packages.html}.

\bibitem[\citeproctext]{ref-webRdocs}
---------. 2026b. \emph{{webR}: {R} in the Browser}. {r-wasm}.
\url{https://docs.r-wasm.org/webr/latest/}.

\bibitem[\citeproctext]{ref-stoudt2021introprelabs}
Stoudt, Sara. 2021. {``Using Learnr Tutorials in an Intro Stats Class as
Pre-Labs.''} Blog post.
\url{https://sastoudt.github.io/posts/2021-06-05-learnr-tutorials-intro-stat/}.

\bibitem[\citeproctext]{ref-Stoudt.Scotina.Luebke:2022}
Stoudt, Sara, Scotina Anthony D., and Karsten Luebke. 2022a.
{``Supporting Statistics and Data Science Education with
\texttt{learnr}.''} \emph{Technology Innovations in Statistics
Education} 14 (1). \url{https://doi.org/10.5070/t514156264}.

\bibitem[\citeproctext]{ref-stoudt2022learnr}
Stoudt, Sara, Anthony D. Scotina, and Karsten Luebke. 2022b.
{``Supporting Statistics and Data Science Education with Learnr.''}
\emph{Technology Innovations in Statistics Education} 14 (1).
\url{https://doi.org/10.5070/T514156264}.

\bibitem[\citeproctext]{ref-Timbers.Campbell.Lee:2022}
Timbers, T., T. Campbell, and M. Lee. 2022. \emph{Data Science: A First
Introduction}. Chapman \& Hall/CRC Data Science Series. CRC Press.
\url{https://www.routledge.com/Data-Science-A-First-Introduction/Timbers-Campbell-Lee/p/book/9780367524685}.

\bibitem[\citeproctext]{ref-ward2026stanplayground}
Ward, Brian, Jeff Soules, and Jeremy Magland. 2026.
{``{stan-playground}: Run Stan Models Directly in Your Browser.''}
\emph{Journal of Open Source Software} 11 (118): 9531.
\url{https://doi.org/10.21105/joss.09531}.

\bibitem[\citeproctext]{ref-white2023webrquarto}
White, Ethan P. 2023. {``A Brief Introduction to Using WebR and Quarto
for Client-Side Interactive Lesson Material.''} Blog post.
\url{https://jabberwocky.weecology.org/2023/03/13/a-brief-introduction-to-using-webr-and-quarto-for-client-side-interactive-lesson-material/}.

\bibitem[\citeproctext]{ref-Woodard.Lee.Woodard:2020}
Woodard, Victoria, Hollylynne Lee, and Roger Woodard. 2020. {``Writing
Assignments to Assess Statistical Thinking.''} \emph{Journal of
Statistics Education} 28 (1): 32--44.
\url{https://doi.org/10.1080/10691898.2019.1696257}.

\bibitem[\citeproctext]{ref-Zeileis.Umlauf.Leisch:2014}
Zeileis, Achim, Nikolaus Umlauf, and Friedrich Leisch. 2014. {``Flexible
Generation of e-Learning Exams in {R}: {M}oodle Quizzes, {OLAT}
Assessments, and Beyond.''} \emph{Journal of Statistical Software} 58
(1): 1--36. \url{https://doi.org/10.18637/jss.v058.i01}.

\bibitem[\citeproctext]{ref-Zimmerman:2020}
Zimmerman, Jane K. 2020. {``Implementing Standards-Based Grading in
Large Courses Across Multiple Sections.''} \emph{PRIMUS\,: Problems,
Resources, and Issues in Mathematics Undergraduate Studies} 30 (8--10):
1040--53. \url{https://doi.org/10.1080/10511970.2020.1733149}.

\end{CSLReferences}




\end{document}
